\section{Columnar tables and vector aggregates}
\label{sec:appendix-columnar}

Columnar storage makes three separate choices visible in a query plan:
which columns to read, which groups of rows to skip, and whether an
aggregate can process a batch of values at once. A query can benefit from
the first two without using a vector aggregate. The physical layout is
described in Section~\ref{sec:pg-storage-columnar}.

\subsection{Create a table with suitable storage options}

The parent table stores scalar columns in stripes. A stripe contains
smaller chunk groups, with separate presence and value streams for each
column. Presence records whether a value is NULL. Minimum and maximum
values can rule out a chunk group before its values are read. Compression
applies to the column streams independently.

Sequence symbols occupy separate storage shards. The parent row identifies
the sequence without requiring its whole payload to be read. Each sequence
column has its own shards. Two sequence columns in one row therefore do
not share a payload merely because their alphabets match.
\texttt{bioseq.len} reports sequence symbols. \texttt{bioseq.unit\_len}
reports packed units, which are not necessarily bytes.

Set initial table options before \texttt{CREATE TABLE}. For example:

\par\noindent\begin{minipage}{\linewidth}
\begin{lstlisting}[style=bioseq-appendix-sql,
  caption={Create a columnar table with two sequence columns.}]
SET bioseq.storage.stripe.row_limit = 20000;
SET bioseq.storage.chunk_group.row_limit = 2000;
SET bioseq.storage.compression.algorithm = 'pglz';

CREATE TABLE run_measurements (
    read_id integer PRIMARY KEY,
    sample_id text NOT NULL,
    read_count integer NOT NULL,
    bases bigint NOT NULL,
    collected_on date NOT NULL,
    forward_read bioseq.seq('dna4', 32) NOT NULL,
    reverse_read bioseq.seq('dna4', 32)
) USING bioseq;

RESET bioseq.storage.stripe.row_limit;
RESET bioseq.storage.chunk_group.row_limit;
RESET bioseq.storage.compression.algorithm;
\end{lstlisting}
\end{minipage}

The relevant settings have the following roles:
\begin{itemize}
  \item \path{bioseq.storage.stripe.row_limit} defaults to 150000 rows.
    \path{bioseq.storage.chunk_group.row_limit} defaults to 10000 rows.
    Each accepts values from 1000 to 100000000.
  \item \path{bioseq.storage.compression.algorithm} accepts
    \texttt{none} and \texttt{pglz}. Builds with the corresponding libraries
    also accept \texttt{lz4} and \texttt{zstd}. The default prefers
    Zstandard, then LZ4, then pglz, according to build support.
    A stream can remain uncompressed when compression does not help.
  \item \path{bioseq.storage.compression.level} defaults to 3 and accepts
    1--19. It affects Zstandard. The current pglz and LZ4 paths do not use
    this level.
  \item \path{bioseq.storage.sequence_shard.max_size_gib} defaults to
    1024 GiB and requires superuser privileges to change. It controls
    later shard allocation. It does not rebalance existing shards.
    Zero starts a new shard for each inserted sequence.
\end{itemize}

The stripe, chunk-group, and compression settings are copied into the
table's options when the table is created. A later \texttt{SET} does not
replace those stored options or rewrite existing data.

\subsection{Change the options of an existing table}

Use the table-option functions to change later writes to an existing table:

\par\noindent\begin{minipage}{\linewidth}
\begin{lstlisting}[style=bioseq-appendix-sql,
  caption={Change stored options, then reset them to current GUC values.}]
SELECT pg_catalog.alter_bioseq_table_set(
    'run_measurements'::regclass,
    stripe_row_limit => 30000,
    chunk_group_row_limit => 3000,
    compression => 'pglz');

SELECT pg_catalog.alter_bioseq_table_reset(
    'run_measurements'::regclass,
    stripe_row_limit => true,
    chunk_group_row_limit => true,
    compression => true,
    compression_level => true);
\end{lstlisting}
\end{minipage}

Both functions require ownership of the table and take an
\texttt{AccessExclusiveLock}. Omitted arguments leave their options
unchanged. The reset function copies the session's current GUC values,
rather than necessarily restoring the original defaults. Neither function
recompresses old stripes.

\subsection{Project, filter, and cache only the needed data}

The following query needs two scalar columns and no sequence symbols:

\par\noindent\begin{minipage}{\linewidth}
\begin{lstlisting}[style=bioseq-appendix-sql,
  caption={Inspect column projection and chunk-group filtering.}]
EXPLAIN (ANALYZE, VERBOSE, COSTS OFF)
SELECT read_id, read_count
FROM run_measurements
WHERE read_count >= 1000;
\end{lstlisting}
\end{minipage}

Read these plan properties separately:
\begin{itemize}
  \item \texttt{Projected Columns} lists columns required by the scan.
  \item \texttt{Chunk Group Filters} describes tests against group
    metadata. \texttt{Groups Removed by Filter}, reported during analyzed
    execution, counts groups that were skipped.
  \item \texttt{Vectorized Filter} describes batch filtering. It does not
    imply that a later aggregate is vectorized.
\end{itemize}

Casting a sequence to \texttt{text} requests its symbols. A slice limits
the requested interval, as described in the sequence-value reference.
The decompressed-column cache can reuse a compressed column buffer during
one scan:

\par\noindent\begin{minipage}{\linewidth}
\begin{lstlisting}[style=bioseq-appendix-sql,
  caption={Enable the decompressed-column cache for one transaction.}]
BEGIN;
SET LOCAL bioseq.column_cache.enabled = on;
SET LOCAL bioseq.column_cache.capacity_mb = 200;

EXPLAIN (ANALYZE, VERBOSE, COSTS OFF)
SELECT read_id, forward_read[0:64]::text
FROM run_measurements
WHERE sample_id = 'alpha';
COMMIT;
\end{lstlisting}
\end{minipage}

The cache is disabled by default. Its capacity defaults to 200 MiB and
accepts 20--20000 MiB. Entries identify a relation, stripe, chunk group,
and column. They belong to one PostgreSQL backend and are cleared when
the scan ends. Plan counters report hits, misses, writes, evictions, and
peak and final cache use. These buffers are separate from the
\texttt{bioseq\_cache} extension, which stores completed query results.
Section~\ref{sec:pg-storage-caches} compares the cache lifetimes.

\subsection{Use a supported vector aggregate}

The current aggregate registry supports:
\begin{itemize}
  \item \texttt{count(*)}, and \texttt{count(column)} for a direct column.
  \item \texttt{sum}, \texttt{avg}, \texttt{min}, and \texttt{max} for
    \texttt{smallint}, \texttt{integer}, and \texttt{bigint} inputs.
  \item \texttt{min} and \texttt{max} for \texttt{date} inputs.
\end{itemize}

An eligible plan has a plain aggregate directly above a
\texttt{BioseqScan}, with \path{bioseq.vectorization.enabled} on.
The following nonempty example has no NULL aggregate inputs. The nullable
second sequence column still illustrates independent sequence storage.

\par\noindent\begin{minipage}{\linewidth}
\begin{lstlisting}[style=bioseq-appendix-sql,
  caption={Load a small table for inspecting the aggregate path.}]
INSERT INTO run_measurements VALUES
    (1, 'alpha', 12, 32, DATE '2026-01-02',
     repeat('ACGT', 8)::bioseq.seq('dna4', 32),
     'TGCATGCATGCATGCA'),
    (2, 'alpha', 20, 32, DATE '2026-01-03',
     'TCACTGTAGGATCGTAGCTAGTCCGTACTGTA', 'ACGTACGTACGTACGT'),
    (3, 'beta', 28, 32, DATE '2026-01-04',
     'GGGGCCCCTTTTAAAAGGGGCCCCTTTTAAAA', NULL),
    (4, 'beta', 36, 32, DATE '2026-01-05',
     'CTAGTCACTGTAGCCTGATCGTAGCTAGTCAA', 'AAAACCCCGGGGTTTT');
ANALYZE run_measurements;
\end{lstlisting}
\end{minipage}

\par\noindent\begin{minipage}{\linewidth}
\begin{lstlisting}[style=bioseq-appendix-sql,
  caption={Ask for and inspect a plain vector aggregate.}]
BEGIN;
SET LOCAL max_parallel_workers_per_gather = 0;
SET LOCAL bioseq.execution.parallel.enabled = off;
SET LOCAL bioseq.vectorization.enabled = on;

EXPLAIN (ANALYZE, VERBOSE, COSTS OFF)
SELECT count(*) AS rows,
       sum(read_count) AS reads,
       avg(bases) AS average_bases,
       min(collected_on) AS first_day,
       max(collected_on) AS last_day
FROM run_measurements;
COMMIT;
\end{lstlisting}
\end{minipage}

The plan to look for is \texttt{Custom Scan (VectorAggNode)} with a
\texttt{BioseqScan} child. Run the \texttt{SELECT} without
\texttt{EXPLAIN} to obtain the values: four rows, a read-count sum of 96,
an average of 32 bases, and dates from 2 to 5 January 2026.
Neither sequence payload is an input to this aggregate.
The two parallel-scan controls disable PostgreSQL query workers for this
example. They do not disable the separate oneTBB task scheduler.

\subsection{Recognize ordinary aggregation and its correctness role}

Grouped aggregates, joins, aggregate \texttt{DISTINCT} or \texttt{FILTER}
clauses, slice projections, and aggregates over an owned
\texttt{mmsearch} result use the ordinary path. Unregistered aggregate
functions and unsupported input expressions do so too. For example:

\par\noindent\begin{minipage}{\linewidth}
\begin{lstlisting}[style=bioseq-appendix-sql,
  caption={Useful queries that do not select the plain vector aggregate.}]
SELECT sample_id, sum(read_count)
FROM run_measurements
GROUP BY sample_id
ORDER BY sample_id;

SELECT max(bioseq.len(forward_read[0:16]))
FROM run_measurements;

SELECT count(DISTINCT sample_id)
FROM run_measurements;
\end{lstlisting}
\end{minipage}

\begin{quote}
\textbf{Current correctness limitation.}
Some vector aggregates return the wrong result for empty or all-NULL
inputs. Small-integer and integer sums can return zero instead of NULL.
Some integer and date minimum or maximum functions can return their
initial extreme value. Keep \path{bioseq.vectorization.enabled} off for
queries that can encounter these cases. The ordinary PostgreSQL aggregate
path provides the expected SQL results.
\end{quote}

\subsection{Inspect and maintain sequence storage}

Use the public shard function to inspect each sequence column:

\par\noindent\begin{minipage}{\linewidth}
\begin{lstlisting}[style=bioseq-appendix-sql,
  caption={Inspect payload shards and their companion relations.}]
SELECT column_attno, bucket_idx, shard_idx,
       sub_relation, sidecar_relation, bytes_written, is_full,
       pg_relation_size(sub_relation) AS payload_relation_bytes,
       pg_relation_size(sidecar_relation) AS sidecar_relation_bytes
FROM bioseq.list_seq_storage_shards(
    'run_measurements'::regclass)
ORDER BY column_attno, bucket_idx, shard_idx;
\end{lstlisting}
\end{minipage}

The column number distinguishes the two sequence columns. The bucket
groups sequences by length, and the shard number distinguishes successive
allocations in that bucket. The two returned relation identifiers name
the payload relation and its companion metadata relation.
\texttt{bytes\_written} and \texttt{is\_full} describe shard allocation.
\texttt{pg\_relation\_size} measures the returned relation's physical size.
Index sizes and update ages are available through the separate index
metadata interface below.

The parent table and its sequence shards have separate maintenance calls:

\par\noindent\begin{minipage}{\linewidth}
\begin{lstlisting}[style=bioseq-appendix-sql,
  caption={Maintain the parent table and allocated sequence shards.}]
VACUUM (ANALYZE) run_measurements;
CALL bioseq.vacuum_seq_storage('run_measurements'::regclass);
\end{lstlisting}
\end{minipage}

Call the procedure outside an explicit transaction block when shards have
been allocated. It runs \texttt{VACUUM ANALYZE} on shard pairs and can
commit between them. An error can leave earlier pairs maintained. After
correcting the error, call it again. A table with no allocated sequence
shards is a no-op. This procedure does not rebuild sequence indexes.

A practical sequence is to choose table options, load rows, run
\texttt{ANALYZE}, and inspect projection and filtering with
\texttt{EXPLAIN ANALYZE}. Enable the column cache or vector aggregation
only when their behavior suits the query. After mutations, maintain the
parent, shards, and indexes according to their separate roles.
Keep \path{bioseq.storage.dml.enabled} on: despite its name, this internal
switch controls reading deletion masks. Turning it off can expose deleted
rows and does not make the table read-only.

\section{Sequence-index maintenance and metadata}
\label{sec:appendix-indexes}

BioSeqDB manages sequence indexes through its own SQL functions.
\texttt{CREATE INDEX} creates a PostgreSQL index and does not build these
sequence search structures. Section~\ref{sec:index-logical-layouts}
explains the layouts. This section gives the operational steps.

\subsection{Choose a family and its update policy}

Each sequence column registers three production index families:
\begin{itemize}
  \item \texttt{minimizers\_per\_fragment} (MPF) stores minimizers for
    sequence fragments. It is the default read family and defaults to
    \texttt{UPDATE\_ON\_INSERT}.
  \item \texttt{minimizers\_per\_sequence} (MPS) groups minimizer
    positions across a sequence. It defaults to
    \texttt{UPDATE\_MANUALLY}.
  \item \texttt{symbol\_counts} stores letter counts and defaults to
    \texttt{UPDATE\_ON\_INSERT}. It serves operations such as
    \texttt{bioseq.symbol\_count(sequence, 'A')}.
\end{itemize}

Registration describes the available families. An empty table has no
sequence-index payload to read. A table with multiple sequence columns
has separate registrations for each column. Both minimizer families can
report \texttt{CHUNKED\_PER\_FRAGMENT} in their physical
\texttt{index\_type}. Use \texttt{index\_name} to distinguish their
logical roles.

There are two levels of insertion maintenance:
\begin{enumerate}
  \item \path{bioseq.index.insert_maintenance.enabled}, on by default,
    permits insertion-time maintenance in the current session.
  \item Each index's stored policy decides whether that family is
    maintained on insertion. \texttt{UPDATE\_MANUALLY} requires an
    explicit rebuild. \texttt{UPDATE\_ON\_INSERT} permits eligible
    insertion-time maintenance.
\end{enumerate}

For an existing \texttt{indexed\_reads} table, the following calls switch
MPS between the two policies:

\par\noindent\begin{minipage}{\linewidth}
\begin{lstlisting}[style=bioseq-appendix-sql,
  caption={Change one index family's policy for every sequence column.}]
SELECT bioseq.set_index_policy(
    'indexed_reads'::regclass,
    'minimizers_per_sequence', 'UPDATE_ON_INSERT');

SELECT bioseq.set_index_policy(
    'indexed_reads'::regclass,
    'minimizers_per_sequence', 'UPDATE_MANUALLY');
\end{lstlisting}
\end{minipage}

Use these exact policy names. The current function treats a string other
than \texttt{UPDATE\_ON\_INSERT} as manual, so a spelling error can
silently select manual maintenance. Changing a policy does not build
missing payloads, clear old payloads, or select the active read family.
The optional grouped construction path has additional restrictions,
described below.

For frequent small insertions, ordinary MPF and symbol counts provide a
useful starting point. For a large load, disable insertion maintenance,
then build the selected families once. For stable data, compare MPF and
MPS using the actual query workload. Letter-composition queries need
symbol counts rather than a minimizer index. A one-off search can use
transient preparation without constructing a persistent index.

\subsection{Load first, then build the selected indexes}

The following example uses a fresh table and commits its source rows
before explicit index construction. The repeated flanks make each target
longer than the query shown later.

\par\noindent\begin{minipage}{\linewidth}
\begin{lstlisting}[style=bioseq-appendix-sql,
  caption={Load a searchable table without insertion-time indexing.}]
BEGIN;
SET LOCAL bioseq.index.insert_maintenance.enabled = off;

CREATE TABLE indexed_reads (
    read_id integer PRIMARY KEY,
    sample_id text NOT NULL,
    sequence bioseq.seq('dna4', 64) NOT NULL
) USING bioseq;

INSERT INTO indexed_reads VALUES
    (1, 'alpha',
     (repeat('ACGT', 96) || 'CTAGTCACTGTAGCCTGATCGTAGCTAGT' ||
      repeat('TGCA', 32))::bioseq.seq('dna4', 64)),
    (2, 'alpha',
     (repeat('TGCA', 72) || 'CTAGTCACTGTAGCCTGATCGTAGCTAGT' ||
      repeat('ACGT', 48))::bioseq.seq('dna4', 64)),
    (3, 'beta', repeat('GATTACA', 80)::bioseq.seq('dna4', 64));
COMMIT;
\end{lstlisting}
\end{minipage}

Inspect the initial state before building:

\par\noindent\begin{minipage}{\linewidth}
\begin{lstlisting}[style=bioseq-appendix-sql,
  caption={Read the public index metadata.}]
SELECT column_name, index_name, index_type, update_policy,
       storage_size, time_since_update
FROM bioseq.list_metadata()
WHERE table_name = 'indexed_reads'::regclass
ORDER BY column_attno, index_name;
\end{lstlisting}
\end{minipage}

There should be three registered rows for this sequence column, with no
built payload and no update time. \texttt{column\_attno} and
\texttt{column\_name} identify the source column.
\texttt{storage\_size} and \texttt{time\_since\_update} describe stored
index data and its last recorded update. They do not establish that a
particular query can use it.

For ordinary relation-backed indexes, leave
\path{bioseq.aligner.grouped_soa_generation.enabled} off and explicitly
build the desired families:

\par\noindent\begin{minipage}{\linewidth}
\begin{lstlisting}[style=bioseq-appendix-sql,
  caption={Build all three ordinary families and select MPF for reads.},
  label={lst:appendix-ordinary-index-build}]
SET bioseq.aligner.grouped_soa_generation.enabled = off;
SET bioseq.aligner.vulkan.enabled = off;
SET bioseq.aligner.minimizer.algorithm = 'two_stack_nthash';
SET bioseq.aligner.minimizer.kmer_size = 15;
SET bioseq.aligner.minimizer.window_size = 10;

SELECT bioseq.update_index(
    'indexed_reads'::regclass, 'minimizers_per_fragment');
SELECT bioseq.update_index(
    'indexed_reads'::regclass, 'minimizers_per_sequence');
SELECT bioseq.update_index(
    'indexed_reads'::regclass, 'symbol_counts');
SELECT bioseq.set_active_read_index(
    'indexed_reads'::regclass, 'minimizers_per_fragment');
ANALYZE indexed_reads;
\end{lstlisting}
\end{minipage}

\texttt{update\_index} returns the number of sequence values processed,
across the table's sequence columns. It is not a count of physical index
rows. \texttt{set\_active\_read\_index} changes the selected family for
all sequence columns. It does not build that family.
The chosen minimizer implementation, $k=15$, and $w=10$ define this
example's sketch profile. Keep the search profile compatible with the
construction profile. Turning Vulkan off selects CPU availability for
this example, rather than rebuilding stored data.

\texttt{ANALYZE} updates PostgreSQL planner statistics.
\texttt{VACUUM}, \texttt{CREATE INDEX}, and \texttt{REINDEX} have their
PostgreSQL roles. None replaces \texttt{bioseq.update\_index}.

\subsection{Inspect physical rows and letter-count reuse}

After construction, inspect the stored index relations through public SQL:

\par\noindent\begin{minipage}{\linewidth}
\begin{lstlisting}[style=bioseq-appendix-sql,
  caption={Inspect index storage and one source row's published keys.}]
SELECT index_name, update_policy, storage_size, time_since_update,
       bioseq.index_table_row_count(
           'indexed_reads'::regclass, index_name) AS physical_rows
FROM bioseq.list_metadata()
WHERE table_name = 'indexed_reads'::regclass
ORDER BY column_attno, index_name;

SELECT r.read_id,
       bioseq.index_distinct_sub_keys(
           'indexed_reads'::regclass,
           'minimizers_per_fragment',
           bioseq.column_row_number(
               'indexed_reads'::regclass, r.ctid)) AS published_keys
FROM indexed_reads AS r
ORDER BY r.read_id;
\end{lstlisting}
\end{minipage}

\texttt{index\_table\_row\_count} counts physical rows across the chosen
family's index shards. \texttt{index\_distinct\_sub\_keys} reports the
largest per-shard count of published sub-keys for the supplied source row
number. A sub-key identifies a stored part of that row's index.
\texttt{column\_row\_number} decodes the current BioSeqDB tuple locator
from \texttt{ctid}. Keep the application key \texttt{read\_id} for durable
row references. A physical locator is not an application identifier.

The letter-count use case is independent of minimizer-family selection:

\par\noindent\begin{minipage}{\linewidth}
\begin{lstlisting}[style=bioseq-appendix-sql,
  caption={Inspect direct use of saved letter counts.}]
EXPLAIN (ANALYZE, VERBOSE, COSTS OFF)
SELECT read_id, bioseq.symbol_count(sequence, 'A') AS adenines
FROM indexed_reads
ORDER BY read_id;
\end{lstlisting}
\end{minipage}

Look for \texttt{Bioseq Fetched From Meta} in a compatible projection
plan. A missing compatible count can be recomputed from the sequence.
Using \texttt{symbol\_count} in a predicate does not create a conventional
PostgreSQL index scan over letter counts.

\subsection{Distinguish permitted reads from required reads}

\path{bioseq.index.sequence_read.enabled}, on by default, permits stored
index reads for ordinary searches. Turning it off prevents such reuse.
It does not delete index data, disable insertion maintenance, or disable
background persistence.

The per-call \texttt{aligner\_options(index\_read\_requirement => ...)}
argument makes a stronger choice:
\begin{itemize}
  \item \texttt{0}: allow a compatible stored index or an eligible
    fallback. Use this for searches that may recompute missing data.
  \item \texttt{1}: force transient target preparation. Existing stored
    indexes may remain present, but the query does not read them.
  \item \texttt{2}: require a compatible stored target index. A missing
    or incompatible index is an error. This overrides the ordinary
    stored-read GUC.
\end{itemize}

These are per-call options, not GUC values. Their relationship to other
search defaults is described in Section~\ref{sec:appendix-search-options}.
A positive physical row count alone does not satisfy requirement 2.
Compatibility also depends on source identity, alphabet, packed-unit
width, minimizer scheme and parameters, stored format, and the chosen
reader. See Section~\ref{sec:index-compatibility}.

For a controlled transient comparison, use requirement 1. To compare the
stored-read GUC on and off, keep requirement 0. Otherwise the per-call
requirement can determine the outcome instead of the GUC.

\subsection{Build and check the experimental grouped reader}

The optional grouped path is a separate construction choice. To try it,
repeat the fresh-table load above under a new table name, or use
\texttt{indexed\_reads} before running
Listing~\ref{lst:appendix-ordinary-index-build}. It should have committed
source rows and no previously built ordinary minimizer layout.

\par\noindent\begin{minipage}{\linewidth}
\begin{lstlisting}[style=bioseq-appendix-sql,
  caption={Build grouped generations from committed source rows.}]
SET max_parallel_workers_per_gather = 0;
SET bioseq.execution.parallel.enabled = off;
SET bioseq.aligner.vulkan.enabled = off;
SET bioseq.aligner.minimizer.algorithm = 'two_stack_nthash';
SET bioseq.aligner.minimizer.kmer_size = 15;
SET bioseq.aligner.minimizer.window_size = 10;
SET bioseq.aligner.chaining.algorithm = 'simd_minimap2';
SET bioseq.aligner.base_alignment.algorithm = 'block_aligner';
SET bioseq.index.async_persistence.enabled = off;
SET bioseq.index.recompute_async_persistence.enabled = off;
SET bioseq.aligner.grouped_soa_generation.enabled = on;

SELECT bioseq.update_index(
    'indexed_reads'::regclass, 'minimizers_per_fragment');
SELECT bioseq.update_index(
    'indexed_reads'::regclass, 'minimizers_per_sequence');
SELECT bioseq.update_index(
    'indexed_reads'::regclass, 'symbol_counts');
SELECT bioseq.set_active_read_index(
    'indexed_reads'::regclass, 'minimizers_per_fragment');
ANALYZE indexed_reads;
\end{lstlisting}
\end{minipage}

The settings in this example have distinct jobs:
\begin{itemize}
  \item \path{bioseq.aligner.grouped_soa_generation.enabled}, off by
    default, changes explicit index construction. Grouped MPF writes
    relation rows and a CPU reader image. Grouped MPS writes relation
    streams. The flag neither builds an index by itself nor selects the
    later query reader.
  \item Both asynchronous-persistence settings must be off for explicit
    grouped construction. Their ordinary purpose is explained below.
  \item The current grouped builder uses $k=15,w=10$. The selected
    minimizer implementation must agree between build and search.
    \texttt{simd\_minimap2} selects CPU chaining, and
    \texttt{block\_aligner} selects base alignment.
  \item Disabling PostgreSQL parallel query workers and Vulkan gives a
    CPU example. BioSeqDB can still schedule its own CPU tasks.
\end{itemize}

Ordinary per-tuple insertion with grouped construction enabled rejects
minimizer families set to \texttt{UPDATE\_ON\_INSERT}. Load with insertion
maintenance off and publish after the source rows commit. The streamed
loader has a separate post-flush batch path, so its maintenance behavior
should not be inferred from ordinary per-tuple insertion.

The current strict reader belongs to an eligible \texttt{BioseqBatchJoin}.
Create a small query table and search its whole sequence column against
the whole target column. A constant-query \texttt{BioseqScan} does not
support this strict reader and rejects requirement 2.

\par\noindent\begin{minipage}{\linewidth}
\begin{lstlisting}[style=bioseq-appendix-sql,
  caption={Prepare query rows for the strict batch-join reader.}]
BEGIN;
SET LOCAL bioseq.index.insert_maintenance.enabled = off;
CREATE TABLE index_queries (
    query_id integer PRIMARY KEY,
    sequence bioseq.seq('dna4', 64) NOT NULL
) USING bioseq;
INSERT INTO index_queries VALUES
    (1, 'CTAGTCACTGTAGCCTGATCGTAGCTAGT');
COMMIT;
ANALYZE index_queries;
\end{lstlisting}
\end{minipage}

\par\noindent\begin{minipage}{\linewidth}
\begin{lstlisting}[style=bioseq-appendix-sql,
  caption={Check stored-index use through an analyzed search.}]
SET bioseq.index.sequence_read.enabled = on;

EXPLAIN (ANALYZE, VERBOSE, COSTS OFF)
SELECT r.read_id, q.query_id,
       bioseq.target_start(r.read_id) AS target_start,
       bioseq.target_end(r.read_id) AS target_end
FROM indexed_reads AS r
JOIN index_queries AS q
  ON mmsearch(r.sequence, q.sequence,
              aligner_options(index_read_requirement => 2))
ORDER BY r.read_id, q.query_id, target_start;

SELECT requested_target_index_source, actual_target_index_source,
       persisted_generation_id, persisted_source_identity,
       persisted_index_rows, missing_persisted_index_rows,
       recomputed_target_index_rows,
       persisted_generation_read_sessions
FROM bioseq.last_alignment_execution_stats()
WHERE stage = 'query_sketch';
\end{lstlisting}
\end{minipage}

For a successful strict check, the analyzed plan should contain
\texttt{BioseqBatchJoin} and report
\texttt{require\_compatible\_persisted} and actual target source
\texttt{persisted}. Missing and recomputed target-index row counts should
be zero. The last-run statistics identify the accepted generation and
source. Read them immediately after the search in the same backend.
They are diagnostic state, not a permanent history of queries.

The grouped MPF image is used by eligible prepared CPU batch-join readers
over whole, untransformed targets. Enabling the
construction flag does not make every search map that file. Ordinary
allow-mode readers can instead read relation-backed payloads. The reader
selection and its limits are described in
Section~\ref{sec:index-example-settings} and
Figure~\ref{fig:index-join-read-paths}.
This path remains experimental. Cleanup of retired images and complete
recovery support are unfinished. Do not treat the example as a durability
guarantee for those images.

Plain \texttt{SET} persists in the session after a successful commit.
After this experiment, restore the changed settings with \texttt{RESET},
or use a dedicated connection. In particular, restore grouped construction
before returning to ordinary insertion-time maintenance.

\subsection{Handle mutations and concurrent use}

Index maintenance follows the identity of the visible source sequence:
\begin{itemize}
  \item \texttt{INSERT} adds a new source value. Eligible indexes follow
    their insertion policy when the session permits maintenance.
  \item \texttt{UPDATE} retires the old row version and inserts a new
    source value. The new value needs index data for its new identity.
  \item \texttt{DELETE} removes the row from visible search input.
    Retiring its source identity is separate from immediately reclaiming
    every stored index chunk.
\end{itemize}

A manually maintained index can retain physical data for an older source
without being suitable for the new one. An allow-mode search may recompute
eligible missing data. A strict search should fail when no compatible
generation is available. Rebuild the selected families after a bulk
change, then run \texttt{ANALYZE}.

PostgreSQL snapshots determine which row versions a scan can see. Index
generation checks must agree with those visible sources. Committing new
source rows and building their indexes in a later transaction creates an
interval during which strict readers can fail. Plan that interval for
maintenance, or allow recomputation where supported. Do not assume that
an asynchronous rebuild becomes available at the source commit.
Section~\ref{sec:pg-storage-versions} explains row visibility, and
Section~\ref{sec:index-generations} explains index publication and its
current durability limits.

Use three checks before depending on stored reads:
\begin{enumerate}
  \item Inspect the per-column family, update policy, size, and update
    time in \nolinkurl{bioseq.list_metadata()}.
  \item Check physical index rows through the public inspection
    functions when storage state needs investigation.
  \item Execute a representative strict search and inspect its actual
    reader choice. Metadata alone does not prove query compatibility.
\end{enumerate}

\subsection{Bound index memory and inspect deferred work}

\path{bioseq.index.memory_cap_mb} limits accounted resident-index
reservations. Zero uses the runtime memory policy. A nonzero value does
not set a whole-query or cluster-wide memory ceiling. Large index builds
can spill sorted runs to PostgreSQL temporary storage and merge them.
Inspect the most recent build and current scheduler state separately:

\par\noindent\begin{minipage}{\linewidth}
\begin{lstlisting}[style=bioseq-appendix-sql,
  caption={Inspect index construction and scheduler reservations.}]
SELECT bioseq.last_index_merge_run_count();

SELECT current_threads, batches_awaiting,
       async_persist_jobs_waiting, index_memory_reserved_bytes
FROM bioseq.scheduler_stats();
\end{lstlisting}
\end{minipage}

The merge-run count is backend-local state from the most recent index
build. It is not a property of every index or a search-time measurement.
The scheduler snapshot reports current threads, waiting batches, queued
persistence jobs, and reserved index bytes. Reserved bytes do not include
every allocation in the process.

Two GUCs control eligible deferred persistence:
\begin{itemize}
  \item \path{bioseq.index.async_persistence.enabled}, on by default,
    permits background persistence requests. It does not turn an explicit
    \texttt{update\_index} call into an asynchronous call.
  \item \path{bioseq.index.recompute_async_persistence.enabled}, off by
    default, also permits an eligible query-time recomputation to request
    persistence for later searches. It requires the preceding setting.
\end{itemize}

Current recomputation persistence needs a real source relation, an
eligible insertion-maintained policy, a supported minimizer scheme, and
$k=15,w=10$. The older deferred path excludes bounded ntHash and does not
build the grouped MPF reader image. Requests are staged for submission
after commit. Abort discards them. A background worker reopens the source
and rechecks eligibility. Missing queue resources or failed eligibility
checks can prevent persistence, so a successful search does not imply
that a stored index was created.

For predictable maintenance, choose the family and profile, choose
insertion or explicit updates, build after loading, and check a real
stored read. After mutations, repeat the relevant build and readiness
steps. Keep source maintenance, PostgreSQL statistics, index publication,
and background work as separate operations when interpreting diagnostics.
